\scanpage{26}
\[
Q^i\subset Q^j\qquad\xi
\]
\[
\deleted{\mathbb Z[T][T']/T^{2\nu+2},\ 2T'-T^{\nu+1}}
\]
\[
\left\{\begin{array}{l}
\xi^0,\ldots,\xi^{\nu-1},[\xi^\nu,\lambda_\nu],\xi_{\nu+1},\ldots,\xi_{2\nu}\\
\xi^0,\ldots,\xi^\nu,\xi_{\nu+1},\ldots,\xi_{2\nu+1}
\end{array}\right.
\]
\margin{\[
\left[\begin{array}{l}
2\xi_i=\xi^i\\
\xi^i_{\deleted{\unsure{i}}}\xi_j=\xi_{i+j}\\
\xi_i\xi_j=0\quad\text{pour raisons de degré}\\
\lambda_\nu^2=2\xi_{2\nu}=\xi^{2\nu}\\
\lambda_\nu\xi=0\\
\lambda_\nu\xi_\alpha=0\quad\alpha\geq\nu+1\\
\text{raisons de degré}
\end{array}\right.
\]}
\diagram{\begin{array}{ccc}
Q^i&\subset&Q^j\qquad(i\leq j)\\
\xi^\alpha&\underset{\deleted{\unsure{\xi}}}{\longleftarrow}&\xi^\alpha\\
\xi_\alpha&\leftsquigarrow&\xi_\alpha\\
0&\leftsquigarrow&\lambda_\nu\\
\xi_\nu&\leftsquigarrow&\varphi_\nu
\end{array}
\quad
\left.\begin{array}{c}\\ \\ \\ \\ \end{array}\right\}\text{si }i<j=2\nu}
\margin{: $\lambda_\nu$ est une \uline{classe \unsure{effaçable}}}
\diagram{\begin{array}{ccc}
\xi^\alpha&\rightsquigarrow&\xi^{\alpha+(j-i)}\\
\xi_\alpha&\rightsquigarrow&\xi_{\alpha+(j-i)}\\
\lambda_\nu&\rightsquigarrow&0\\
\varphi_\nu&\rightsquigarrow&\xi_{\nu+(j-i)}
\end{array}\quad
\left.\begin{array}{c}\\ \\ \\ \end{array}\right\}\text{si }i=2\nu<j}
\margin{: $\lambda_\nu$ est une classe ``évanescente''.}
\mspage{25}
